\paragraph{Separating redundancy from the height gap.}
The code size and the geometric height bound need not use the same
exponent. The following estimate makes their different roles explicit.

\begin{proposition}[A linear source-spread cost]\label{shared-code-trace}
Let \(T_i:\mathbb C^{r_i}\to\mathbb C^D\) be a finite family, put
\[
 R=\sum_i r_i>0,\quad Q=\sum_iT_iT_i^*,\quad
 \mathcal S=\bigoplus_iT_i^*T_i,\quad M=\operatorname{tr}Q,
\]
and let \(D_*>0\), \(\xi=D_*/R\), and \(0<E\le1/64\).
Assume \(\xi\le1/64\) and
\[
 \|\mathcal S-I_R\|_{\rm HS}^2\le E^2R,\qquad
 D\le(1+E)D_*,\qquad
 \operatorname{tr}Q^2\le(1+E)\frac{R^2}{D_*}+(1+E)^2R.
\]
Let \(J_i\) be isometric source inclusions whose total discarded
dimension is at most \(E^2R\). Suppose the first \(n\) rows of
\(T_iJ_i\), where \(1\le n\le D\) and \(1-n/D\le E\), represent
the same integral arithmetic maps in one common target basis.
Assume the height and slope hypotheses of
Corollary~\ref{shared-trace-ambient} for these final source lattices
and this target, with \(U,H\ge0\). Then
\begin{equation}\label{shared-code-trace-bound}
 s-\mu(F)\le\min\{H,\ U+(8\sqrt E+2\sqrt\xi)H\}.
\end{equation}
\end{proposition}

\begin{proof}
Cauchy--Schwarz gives \((1-E)R\le M\le(1+E)R\).
Define
\[
 \gamma=\frac{\operatorname{tr}((M/R)I_R-\mathcal S)_+}{M},\quad
 \beta=\frac{\operatorname{tr}(Q-(M/D)I_D)_+}{M},\quad
 \eta=\frac{D\operatorname{tr}Q^2}{M^2}-1 .
\]
Centering lowers the squared source error, and a trace-zero matrix
has positive trace equal to half its trace norm. Therefore
\[
 \gamma\le\frac{E}{2(1-E)}\le\frac{2E}{3},\qquad
 \eta\le\frac{4E+(1+E)^3\xi}{(1-E)^2}\le8E+4\xi,\qquad
 \beta\le\frac{\sqrt\eta}{2}\le2\sqrt E+\sqrt\xi .
\]
For the middle bound, \(4/(1-E)^2\le64/9<8\) and
\((1+E)^3/(1-E)^2\le125/36<4\) already hold for \(E\le1/4\).
The fraction \(\rho\) of energy lost by the source restrictions is
at most \(E^2+\gamma\le E\), by the discarded-projection estimate.
Put \(\delta=1-n/D\) and \(q=\gamma+\beta+\rho\).
Equation~\eqref{shared-trace-source-only} applies and bounds the
uncapped cost by
\[
 \frac{\delta U+q(U+H)/(1-\rho)}{1-\delta}.
\]
Here \(q\le2E+2\sqrt E+\sqrt\xi\). Its coefficient of \(U\) is at most
\((E+q)/(1-E)^2\le(27/64)/(63/64)^2<1\), while its coefficient of
\(H\) is at most \(2q\le8\sqrt E+2\sqrt\xi\).
The independent cap \(H\) proves the assertion.
\end{proof}

\begin{corollary}\label{shared-code-height-separation}
Fix \(\kappa,t>0\), \(h\ge0\), and \(\omega\ge2h+2t\).
Suppose the preceding hypotheses hold at every sufficiently large
scale \(x\), with
\[
 E=e^{-x^\kappa},\qquad \xi\le x^{-\omega},\qquad
 H\le x^h b,\quad b>0.
\]
Then, eventually and uniformly over all the maps covered at that scale,
\[
 s-\mu(F)\le\min\{H,\ U+10b\,x^{-t}\}.
\]
\end{corollary}
\begin{proof}
Eventually \(x^{h+t}\sqrt E\le1\); the two terms in
\eqref{shared-code-trace-bound} then contribute at most
\(8b x^{-t}\) and \(2b x^{-t}\). No growth bound on \(U\) is used.
\end{proof}

The required choices must belong to one construction. For example,
suppose its code has \(\lceil3\Gamma^2x^\omega\rceil\) words,
each full source rank is at least \(D_*/(3\Gamma^2)\), and this code
fits at the same geometric ranks. Then \(R\ge D_*x^\omega\).
Suppose also that the construction and the height estimate have
finitely many staged thresholds: their middle thresholds depend only
on fixed outer data, and their last thresholds depend only on the
already chosen middle entries. Choose \(h\ge1\) above every middle
threshold, then choose \(\omega\ge2h+2t\), and finally choose each
last constant above all required threshold values at both arguments.
No monotonicity of the threshold functions is needed. A constructor
uniform over these choices yields one family satisfying both
requirements; a larger last constant does not reset \(h\).
The global matrix estimates, arithmetic height bound, and compatible
constructor remain substantive hypotheses. The surviving source
spread \(U\) and the other arithmetic losses still require their
own bounds.
